% !TeX root = ../solution.tex

\hypertarget{he22.34}{%
\chapter{[HE22.34] AES Burgers}\label{he22.34}}

\begin{marginfigure}
	\includegraphics[width=49mm]{level8/challenge34.jpg}
\end{marginfigure}
\section{Intro}
Welcome to AES Burgers - where the patty is tatty™ !

\noindent Connect to our server, and place your order!

\noindent \verb+nc 46.101.107.117 2207+

\noindent File: \verb+aesburgers.py+


\section{Solution}\label{hv22.34solution}

The program \verb+aesburgers.py+ supplied does one thing: it asks for a bun and a number of patties, builds a burger and then encrypts the burger with AES in electronic cookbook mode (ECB).  The relevant two functions are:

{\small
\begin{minted}{python}
def encrypt(plain):
    cipher = AES.new(key.encode(), AES.MODE_ECB)
    plain += b' ' * ((16 - (len(plain) % 16)) % 16)
    enc = cipher.encrypt(plain)
    return ''.join('%02x' % b for b in enc)

def makeBurger(bun, patties):
    burger = b''
    burger += bun[::-1] # the bottom bun, flipped
    burger += patties * flag # the patties
    burger += bun # the top bun
    return burger
\end{minted}
}

\noindent
Since the burger is encrypted using ECB, we can craft a known plaintext attack:
In ECB, a given 16-byte-aligned 16-byte-block is always encrypted to the same
result.  If we can construct a bun in such a way that the reversed bun
corresponds to a given block of the burger, then we can brute force the flag.

To illustrate the principle, assume a bun \verb+0123456789ABCDEF+ and assume
that the number of patties times the flag-length is $n\times 16 + 1$.  Then the 
constructed burger looks like
\noindent
\begin{verbatim}
FEDCBA9876543210
he2022{...
...
}0123456789ABCDE
F
\end{verbatim}
\noindent
Using the bun, we have full control over the first 16-bytes of the burger and we know 15
out of the second to last 16-byte block.  The one, unknown byte is determined by the flag. 

If we are now choosing a bun like \verb+AAAAAAAAAAAAAAA.+, then the 
first block becomes \verb+.AAAAAAAAAAAAAAA+, the second last will be
\verb+}AAAAAAAAAAAAAAA+.  This allows us to try all possible buns by 
replacing the \verb+.+ with any printable character and check if the encrypted
blocks match.  If they match, we know the last character of the flag (which we 
assume to be \verb+}+).  Once we have one character, we try to brute force
the next character by changing the number of patties and constructing a new 
bun \verb+AAAAAAAAAAAAAA}.+

The first step is to figure out the number of characters in the flag.  Try all
allowed values for the number of patties, observe the length of the returned, 
encrypted burger and solve for the smallest possible length of the flag.  Can 
be done by thinking or just by brute forcing, and find 35 as flag length.

With this we have all information to start the brute force attack.  Once 16
characters have been found, the third to last block has to be considered and
the bun is just the reversed 16 bytes of the flag considered.

{\small
\begin{minted}{python}
from pwn import *
IP = '46.101.107.117'
PORT = 2207

def patties(flag_len, offset):
    for n in range(1,25):
        if (n*flag_len)%16 == offset%16:
            return n
    print(f'[x] cannot find number of patties')
    exit(-1)


p = remote(IP, PORT)

bun = b'1'*16
flag_len = 35
flag = ''

for j in range(1,flag_len+1):
    n = patties(flag_len, j)
    print(f'[ ] need {n} patties for char {j}')

    for i in range(33,127):
        bun = chr(i) + flag
        if len(bun) < 16:
            bun += (16-j)*'a'
        else:
            bun = bun[:16]
        bun = bun[::-1].encode()
        p.recvuntil(b'How many patties: ')
        p.sendline(bytes(str(n), encoding='ASCII'))

        p.recvuntil(b'Which bun? ')
        p.sendline(bun)
        p.recvuntil(b"Heres your order, enjoy!\n")
        burger = p.recvuntil(b'\n')
        cnub = burger[:32]
        shift = ((j-1)//16)*32
        cbunf = burger[-65-shift:-33-shift]
        if cnub == cbunf:
            flag = chr(i) + flag
            print(f'[+] found {j}th character ', chr(i), ' flag: ', flag)
            break
p.close()
\end{minted}
}
Flag: \verb+he2022{w3_luv_junk_f00d_s0m3t1m3s!}+





	









